\section{Introduction}
\IEEEPARstart{I}{ndustrial-agent} deployments must connect each task to its
production context, process signals, and equipment interfaces. Existing work
addresses end-to-end automation control, operator-reviewed control-sequence
generation, modular factory skills, and production scheduling
\cite{xia2025control,huang2026made,koehle2026llmdriven,waseem2026pretrained}.
Other studies target PLC code generation, context-aware control, resource
capability discovery, shop-floor coordination, digital-twin fault handling,
natural-language machine access, and robotic-line control
\cite{liu2026agents4plc,wu2025instructmpc,lim2025llmenabled,zhao2026shopfloor,gill2025llmdt,hofmann2025nlc}.
Across production lines, agent roles, signal maps, protocols, objectives, and
product/run context vary. An integration architecture must accommodate these
differences around shared task and operation services while keeping each
action tied to its active production context. Recent industrial-agent
architectures add AAS-grounded OPC UA actions with signed approval,
P\&ID-grounded proposal checks, and governed sensor--MES--ERP evidence
streams \cite{graul2026aas,schall2026safe,lin2026agentic}. These works address
asset-specific execution, process validation, or closed-loop analytics. This
paper addresses the complementary software-integration question of composing
task and harness routes, production context, semantic protocol I/O, and
route-owned records. Its simulator evaluation tests configured paths rather
than process-control performance or portability.

Consider two production scenes with different process signals and interfaces:
cast-film thickness control and annealing-zone temperature control. An agent
task in either scene needs production context, typed signal access, and a
route from a named operation to the configured device protocol. The framework
must also distinguish a command receipt from a subsequent process
observation. The shared task and operation services can therefore be
configured with scene-specific signal maps, protocols, and limits while
retaining a common execution path. The evaluation instantiates this path in
two simulator scenes.

This paper presents AgentWorkShop as a reference implementation of
task-to-device integration for industrial agents. It connects collaborative tasks
and registered harnesses to production context, manufacturing execution
system (MES) data, analysis and modeling services, data-acquisition (DAQ) and
write-control (DCW) nodes, and protocol-specific drivers. Its interfaces include
command-line, web, SDK, REST, Model Context Protocol (MCP), and Agent-to-Agent
(A2A) surfaces. Task and chat routes enter the
collaboration runtime, whereas a direct REST tool-invocation route can call
the shared operation dispatcher without creating a task. This direct entry
path is described from the implementation but was not exercised in R20 or
R21 (Table~\ref{tab:coverage}).
Scene blueprints define signal maps, objectives, and guards; channel profiles,
tool policies, and agent bindings scope operations. Agents and callers invoke
named operations through shared runtime checks; the framework does not choose
routes autonomously and retains proposals, command confirmation, and process
observations as distinct records. Fig.~\ref{fig:arch} presents the system
architecture, and Fig.~\ref{fig:runtime} shows the cross-module task flow.
Section~IV details the node-exploration and recipe lifecycles.

This paper asks whether a configured registered-harness path can execute
bounded tasks in two simulator scenes with different signal and protocol
mappings while retaining a checkable relation from task through approval and
verified route receipt to a subsequent run-tagged measurement. A deterministic
fixture separately replays the software route without external model
decisions. A focused mock-state suite changes each line, product, recipe, and
run identifier while a recipe action awaits approval; stale actions are
refused before mock-driver writes. The experiments do not test configuration
reuse, cross-harness portability, or comparative integration benefit.

The paper makes two contributions:
\begin{enumerate}
\item It specifies and implements a modular task-to-device reference path that
  joins collaborative tasks and registered harnesses to turn-time production
  context, per-call operation checks, semantic I/O, configured protocol
  drivers, and service-owned records. For approved recipe actions, the
  implementation binds the approval to a line--product--recipe--run tuple and
  refuses execution if that tuple changes while approval is pending. This
  distinguishes prompt grounding, dispatch authority, device receipt, and
  later process observation; it is an implementation contribution, not a new
  protocol, safety-verification method, or process-control algorithm.
\item It evaluates selected paths in cast-film and annealing simulator scenes.
  In R20, six sequential tasks through one registered harness passed the
  task-to-approval-to-receipt-to-sample linkage checks; none of the six samples
  reached its target band. In R21, a deterministic fixture executed one task
  per scene with no external model and passed 31 runner assertions, including
  cleanup checks. A separate mock-state suite refused each stale
  line--product--recipe--run context before any mock-driver write. These
  checks characterize selected
  software paths; they do not establish process attainment, portability, or
  performance superiority.
\end{enumerate}

Section~II reviews related industrial-agent and integration work.
Section~III presents the design rationale, architecture, and cross-module
execution flow; Section~IV explains module responsibilities and their
integrated operation. Section~V describes the experiments and results, and
Sections~VI and~VII discuss the evidence limits and conclude the paper.
